% ===================================================================
%  Mathemission worksheet template  (Mathemission Teach)
%  Compile with pdfLaTeX in Overleaf, MiKTeX or TeX Live.
%  Optional: put mathemission-logo.png next to this file for the logo
%  (download it from teach.mathemission.in/worksheet/templates/).
%
%  How to use
%   1. Fill in the details under "1 · Sheet details".
%   2. Keep the sections you need (Single, Multiple, Numerical,
%      Matching, Paragraph, Subjective) and delete the rest.
%   3. Start each question with \Q, then the options or answer space:
%         \Q  Question text with $maths$.
%         \OptFour{a}{b}{c}{d}      % short options, 4 across
%         \OptTwo{a}{b}{c}{d}       % medium options, 2 across
%         \OptOne{a}{b}{c}{d}       % long options, one per line
%         \AnsBox                   % box for a numerical answer
%         \Work{4cm}                % space for working (subjective)
%   4. Fill the answer key at the end (or delete it).
% ===================================================================
\documentclass[11pt]{article}
\usepackage[a4paper,top=14mm,bottom=18mm,left=14mm,right=14mm]{geometry}
\usepackage[T1]{fontenc}
\usepackage{lmodern}
\usepackage{amsmath,amssymb}
\usepackage{xcolor,graphicx,tikz}
\usepackage{tabularx,array,multicol}
\usepackage{enumitem}
\usepackage{fancyhdr,lastpage}
\usepackage[most]{tcolorbox}

% ---------- Mathemission colours ----------
\definecolor{mmnavy}{HTML}{0B1F4D}
\definecolor{mmblue}{HTML}{1D4FC4}
\definecolor{mmgold}{HTML}{B7860B}
\definecolor{mmtint}{HTML}{F2F6FD}
\definecolor{mmrule}{HTML}{C9D3EA}
\definecolor{mmgrey}{HTML}{5A6788}

% =================== 1 · Sheet details ===================
\newcommand{\Inst}{Mathemission}                  % your institute / school
\newcommand{\Teacher}{Your name}                  % teacher
\newcommand{\Exam}{JEE Advanced}                  % JEE Advanced / JEE Main / CBSE XII ...
\newcommand{\SheetTitle}{Definite Integration · Practice Sheet 3}
\newcommand{\Batch}{XII-A JEE}
\newcommand{\SheetDate}{15 Nov 2026}
\newcommand{\Duration}{60 minutes}
\newcommand{\MaxMarks}{60}
\newcommand{\SetLetter}{A}                        % leave empty {} for no set badge
% ==========================================================

% ---------- page footer ----------
\pagestyle{fancy}\fancyhf{}
\renewcommand{\headrulewidth}{0pt}
\fancyfoot[L]{\color{mmgrey}\footnotesize\sffamily \Inst\ \textperiodcentered\ \Teacher}
\fancyfoot[R]{\color{mmgrey}\footnotesize\sffamily Page \thepage\ of \pageref{LastPage}}

\setlength{\parindent}{0pt}
\setlength{\parskip}{3pt}

% ---------- header ----------
\newcommand{\Masthead}{%
  \noindent
  \begin{minipage}[c]{14mm}\IfFileExists{mathemission-logo.png}{\includegraphics[width=13mm]{mathemission-logo.png}}{%
    \tikz\node[fill=mmnavy,text=white,rounded corners=2.5mm,minimum size=13mm,font=\bfseries\sffamily\Large]{M};}\end{minipage}\hspace{2mm}%
  \begin{minipage}[c]{\dimexpr\linewidth-38mm}
    {\sffamily\bfseries\LARGE\color{mmnavy}\Inst}\\[1pt]
    {\sffamily\scriptsize\color{mmgrey}\MakeUppercase{\Exam\ \textperiodcentered\ \Teacher}}
  \end{minipage}\hfill
  \ifx\SetLetter\empty\else
  \begin{tikzpicture}[baseline=(b.center)]
    \node[draw=mmnavy,line width=1pt,rounded corners=2mm,inner sep=2mm,align=center,font=\sffamily\bfseries\color{mmnavy}] (b)
      {{\scriptsize SET}\\[-1pt]{\LARGE\SetLetter}};
  \end{tikzpicture}\fi\par\vspace{2mm}
  {\color{mmnavy}\rule{\linewidth}{1.4pt}}\par\nointerlineskip\vspace{0.6mm}
  {\color{mmgold}\rule{46mm}{1.4pt}}\par\vspace{3mm}
  {\sffamily\bfseries\Large \SheetTitle}\par\vspace{2.5mm}
  \begin{tcolorbox}[colback=white,colframe=mmrule,boxrule=0.6pt,arc=1.5mm,left=0mm,right=0mm,top=1mm,bottom=1mm]
    \sffamily\small
    \begin{tabularx}{\linewidth}{@{\hspace{3mm}}X|X|X|X}
      {\tiny\color{mmgrey}\bfseries BATCH}\newline \Batch &
      {\tiny\color{mmgrey}\bfseries DATE}\newline \SheetDate &
      {\tiny\color{mmgrey}\bfseries TIME}\newline \Duration &
      {\tiny\color{mmgrey}\bfseries MAXIMUM MARKS}\newline \MaxMarks
    \end{tabularx}
  \end{tcolorbox}
  \vspace{2mm}
  {\sffamily\small\color{mmgrey}%
   Name \hrulefill\hspace{6mm} Roll no. \rule{30mm}{0.4pt}\hspace{6mm} Marks \rule{20mm}{0.4pt}}\par\vspace{4mm}}

% ---------- instructions panel ----------
\newtcolorbox{instructions}{enhanced,colback=mmtint,colframe=mmtint,boxrule=0pt,arc=0pt,
  borderline west={2.4pt}{0pt}{mmnavy},left=4mm,right=4mm,top=2mm,bottom=2mm,
  fontupper=\sffamily\small,before upper={{\color{mmnavy}\bfseries\scriptsize INSTRUCTIONS}\par\smallskip}}

% ---------- section bar:  \Section{1}{Single correct}{+3 / -1 each} ----------
\newcommand{\Section}[3]{\par\leftskip=0pt\vspace{5mm}\noindent
  \tikz[baseline=(t.base)]\node[fill=mmnavy,text=white,rounded corners=1.2mm,inner xsep=2.5mm,inner ysep=1.2mm,font=\sffamily\bfseries\scriptsize](t){SECTION #1};%
  \hspace{2mm}{\sffamily\bfseries\color{mmnavy}#2}\hspace{2mm}{\color{mmrule}\leaders\hrule height 0.6pt depth -0.2pt\hfill}%
  \hspace{2mm}{\sffamily\small\color{mmgrey}#3}\par\vspace{3mm}}

% ---------- questions ----------
\newcounter{qn}
\newcommand{\Qcircle}[1]{\tikz[baseline=(c.base)]\node[circle,fill=mmnavy,text=white,inner sep=0pt,minimum size=6.2mm,font=\sffamily\bfseries\footnotesize](c){#1};}
\newenvironment{qbody}{\begin{list}{}{\leftmargin=9mm\labelwidth=7mm\labelsep=2mm\itemsep=0pt\topsep=0pt}}{\end{list}}
% Questions sit 9 mm in from the margin; the number hangs in that space.
\newcommand{\Q}{\par\vspace{3.2mm}\stepcounter{qn}\leftskip=9mm\noindent\llap{\makebox[9mm][l]{\Qcircle{\theqn}}}\ignorespaces}
\newcommand{\Marks}[1]{\hfill{\sffamily\bfseries\footnotesize\color{mmnavy}[#1]}}   % e.g. \Marks{4} for subjective
\newcommand{\Opt}[1]{{\sffamily\bfseries\small\color{mmnavy}(#1)}~}
\newcommand{\OptFour}[4]{\par\vspace{1mm}\noindent\begin{tabularx}{\dimexpr\linewidth-\leftskip}{@{}XXXX@{}}\Opt{A}#1&\Opt{B}#2&\Opt{C}#3&\Opt{D}#4\end{tabularx}\par}
\newcommand{\OptTwo}[4]{\par\vspace{1mm}\noindent\begin{tabularx}{\dimexpr\linewidth-\leftskip}{@{}XX@{}}\Opt{A}#1&\Opt{B}#2\\[1mm]\Opt{C}#3&\Opt{D}#4\end{tabularx}\par}
\newcommand{\OptOne}[4]{\par\vspace{1mm}\begin{list}{}{\leftmargin=16mm\labelwidth=7mm\labelsep=0pt\itemsep=1mm\topsep=0pt}
  \item[\Opt{A}\hfill]#1\item[\Opt{B}\hfill]#2\item[\Opt{C}\hfill]#3\item[\Opt{D}\hfill]#4\end{list}}
\newcommand{\AnsBox}{\par\vspace{1.5mm}\noindent{\sffamily\scriptsize\color{mmgrey}ANSWER}\hspace{2mm}\tikz[baseline=1mm]\draw[mmgrey,rounded corners=1mm](0,0)rectangle(32mm,7mm);\par}
\newcommand{\Work}[1]{\par\vspace{1.5mm}\noindent\tikz\draw[mmrule,dashed,rounded corners=1.5mm](0,0)rectangle(\dimexpr\linewidth-\leftskip\relax,#1);\par}
% Paragraph (passage) box, then its questions with \Q as usual
\newtcolorbox{passage}[1][]{colback=mmtint,colframe=mmrule,boxrule=0.6pt,arc=1.5mm,left=3mm,right=3mm,top=2mm,bottom=2mm,
  before upper={{\sffamily\bfseries\scriptsize\color{mmnavy}PARAGRAPH #1}\par\smallskip},before={\par\leftskip=0pt\vspace{2mm}}}
% List-I / List-II for matching questions
\newcommand{\MatchLists}[2]{\par\vspace{1mm}\noindent\begin{tabularx}{\dimexpr\linewidth-\leftskip}{@{}X|X@{}}
  \multicolumn{1}{c|}{\sffamily\bfseries\small List-I}&\multicolumn{1}{c}{\sffamily\bfseries\small List-II}\\\hline #1 & #2\end{tabularx}\par}

% ---------- answer key ----------
\newcommand{\KeyRow}[8]{{\sffamily\bfseries\color{mmnavy}#1}&#2&{\sffamily\bfseries\color{mmnavy}#3}&#4&{\sffamily\bfseries\color{mmnavy}#5}&#6&{\sffamily\bfseries\color{mmnavy}#7}&#8\\\hline}

\begin{document}
\Masthead

\begin{instructions}
\begin{itemize}[leftmargin=4mm,itemsep=1pt,topsep=0pt]
  \item \textbf{Single correct:} only one option is correct. $+3$ if correct, $-1$ if wrong, $0$ if left blank.
  \item \textbf{Multiple correct:} one or more options are correct. $+4$ if all correct options and no wrong option are chosen; $+1$ for each correct option chosen when no wrong option is chosen; $-2$ if any wrong option is chosen.
  \item \textbf{Numerical:} the answer is a number. $+4$ if correct, $0$ otherwise.
  \item \textbf{Matching:} $+3$ if correct, $-1$ if wrong.
  \item \textbf{Subjective:} write the full working; marks are shown against each question.
\end{itemize}
\end{instructions}

% ======================= SECTION 1 =======================
\Section{1}{Single correct}{4 Q \textperiodcentered\ +3 / $-1$ each}

\Q The value of $\displaystyle\int_0^{\pi/2}\frac{\sin x}{\sin x+\cos x}\,dx$ is
\OptFour{$\dfrac{\pi}{2}$}{$\dfrac{\pi}{4}$}{$\pi$}{$0$}

\Q Let $f(x)=x^3-3x+1$. The number of real roots of $f(x)=0$ in $[-2,2]$ is
\OptFour{$0$}{$1$}{$2$}{$3$}

\Q If $\displaystyle\lim_{x\to0}\frac{1-\cos(ax)}{x^2}=8$, then the positive value of $a$ is
\OptTwo{$2$}{$4$}{$8$}{$16$}

\Q Which of the following statements about $g(x)=|x|+|x-1|$ is true?
\OptOne{$g$ is differentiable everywhere}{$g$ has its minimum value $1$ on the whole interval $[0,1]$}{$g$ is strictly increasing on $\mathbb{R}$}{$g$ is unbounded below}

% ======================= SECTION 2 =======================
\Section{2}{Multiple correct}{2 Q \textperiodcentered\ +4 / $-2$ \textperiodcentered\ partial marks}

\Q Let $I_n=\displaystyle\int_0^1 x^n e^{x}\,dx$. Then
\OptOne{$I_n+nI_{n-1}=e$}{$I_0=e-1$}{$I_1=1$}{$I_2=e-2$}

\Q For $z=\cos\theta+i\sin\theta$, which are always true?
\OptTwo{$|z|=1$}{$\bar z=\dfrac1z$}{$z^2=\cos2\theta+i\sin2\theta$}{$z+\bar z=0$}

% ======================= SECTION 3 =======================
\Section{3}{Numerical}{2 Q \textperiodcentered\ +4 / 0 each}

\Q The number of ways of arranging the letters of the word \textsf{ARRANGE} so that the two R's are together is
\AnsBox

\Q The value of $\displaystyle\sum_{r=1}^{10} r^2$ is
\AnsBox

% ======================= SECTION 4 =======================
\Section{4}{Matching}{1 Q \textperiodcentered\ +3 / $-1$}

\Q Match each function in List-I with its range in List-II.
\MatchLists{(P) $\sin^2 x$\newline (Q) $e^{-x^2}$\newline (R) $x^2+1$\newline (S) $\tan^{-1}x$}%
           {(1) $(0,1]$\newline (2) $[0,1]$\newline (3) $[1,\infty)$\newline (4) $\left(-\frac\pi2,\frac\pi2\right)$}
\OptTwo{P$\to$2; Q$\to$1; R$\to$3; S$\to$4}{P$\to$1; Q$\to$2; R$\to$3; S$\to$4}{P$\to$2; Q$\to$1; R$\to$4; S$\to$3}{P$\to$3; Q$\to$1; R$\to$2; S$\to$4}

% ======================= SECTION 5 (paragraph) =======================
\Section{5}{Paragraph}{2 Q \textperiodcentered\ +3 / 0 each}

\begin{passage}[I]
Let $f:\mathbb{R}\to\mathbb{R}$ be defined by $f(x)=x^2-4x+3$.
\end{passage}

\Q The minimum value of $f$ is
\AnsBox

\Q The sum of the roots of $f(x)=0$ is
\AnsBox

% ======================= SECTION 6 (subjective) =======================
\Section{6}{Subjective}{marks as shown}

\Q Prove that $\displaystyle\int_0^{a} f(x)\,dx=\int_0^{a} f(a-x)\,dx$ and hence evaluate $\displaystyle\int_0^{\pi}\frac{x\sin x}{1+\cos^2x}\,dx$. \Marks{5}
\Work{5cm}

\par\leftskip=0pt\vspace{6mm}
\begin{center}{\sffamily\bfseries\footnotesize\color{mmrule} --- END ---}\end{center}

% ======================= ANSWER KEY (new page) =======================
\newpage
\Masthead
{\sffamily\bfseries\color{mmnavy} ANSWER KEY}\par\vspace{2mm}
\begin{tcolorbox}[colback=white,colframe=mmrule,boxrule=0.6pt,arc=1.5mm,left=0mm,right=0mm,top=0mm,bottom=0mm]
\renewcommand{\arraystretch}{1.5}
\begin{tabularx}{\linewidth}{@{\hspace{3mm}}l X | l X | l X | l X}
\KeyRow{1}{(B)}{2}{(D)}{3}{(B)}{4}{(B)}
\KeyRow{5}{(A), (B), (C), (D)}{6}{(A), (B), (C)}{7}{360}{8}{385}
\KeyRow{9}{(A)}{10}{$-1$}{11}{4}{12}{$\pi^2/4$}
\end{tabularx}
\end{tcolorbox}

\end{document}
